\section{Related Work}
\label{sec:related}

\paragraph{Web agent benchmarks.}
WebArena \citep{zhou2024webarena} provides self-hosted replicas of four web
applications with \num{812} long-horizon tasks and a functional evaluator that
scores the resulting environment state rather than the action sequence. Its
reported GPT-4 agent success rate of \SI{14.41}{\percent} against a human
\SI{78.24}{\percent} established the long-horizon gap that motivates this line of
work. VisualWebArena \citep{koh2024visualwebarena} adds visually grounded tasks
and shows that a multimodal agent reaches only \SI{16.4}{\percent} against a human
\SI{88.7}{\percent}. Both benchmark agent \emph{capability} under clean
conditions; neither injects faults, and both report a single success number per
configuration.

\paragraph{Safety and prompt-injection evaluation.}
WASP \citep{evtimov2025wasp} and SecureWebArena \citep{levy2025securewebarena}
treat the page as an adversarial surface, measuring intermediate attack success
separately from end-to-end goal success and introducing process-level rates such
as reasoning-compromise and payload-delivery. Their lesson --- that a
single outcome number cannot describe a multi-step agent's failure modes ---
applies directly to non-adversarial settings, which is where our work sits.

\paragraph{Fault injection for tool-using agents.}
ChaosLLM \citep{zhang2024chaosllm} interposes a middleware between an agent and
its tools to inject unreachability, slowness, and error responses, and finds
hangs the most damaging of the fault types it studies. WAREX
\citep{zhang2025warex} uses a transparent HTTP proxy to inject network faults and
popups into WebArena itself without modifying agent or benchmark, and reports
success-rate drops above \SI{70}{\percent} under network errors. ReliabilityBench
\citep{gupta2025reliabilitybench} combines repeatability, input perturbation, and
infrastructure-fault tolerance into a single reliability surface and compares
ReAct with Reflexion. These systems establish that faults matter; our focus is
different --- we ask what the \emph{measurement} can and cannot resolve, and we
show that the paired design and the evaluator stratum, not the effect size alone,
determine whether a conclusion is warranted.

\paragraph{Multi-agent fault attribution.}
MAS-FIRE \citep{wang2025masfire} catalogues fifteen fault classes across planning,
memory, reasoning, action, and communication, and injects them non-invasively
into MetaGPT, CAMEL, and Table-Critic. It distinguishes fault detection, local
repair, and final task success --- a process-level decomposition we adopt in
spirit but cannot replicate empirically here, because the recovery labels
available in our artifacts are degenerate (\secref{sec:results}).

\paragraph{Gap.}
What is missing is a study that (i) injects non-adversarial faults into WebArena
while keeping the \emph{official} evaluator as the outcome, (ii) uses a paired
design so that the comparison is not swamped by task heterogeneity, and
(iii) reports the power of that design rather than interpreting its nulls as
robustness. We provide exactly that, and the answer we get is methodological: the
outcome measure is the bottleneck.
